% --- PREAMBLE ---
\pagenumbering{roman}

% Title Page
\begin{titlepage}
    \centering
    \vspace*{1cm}
    \Huge \textbf{Capstone Project: Exploring Convolutional Neural Processes for Weather Downscaling} \\
    \vspace{1.5cm}
    \Large \textbf{Student Name:} Francisco Passos \\
    \Large \textbf{Student ID/LEGI:} 24-952-418 \\
    \vspace{1cm}
    \large \textbf{Course:} DAS in Data Science \\
    \large \textbf{Institution:} ETH Zürich \\
    \vspace{2cm}
    \large \textbf{Submission Date:} February 13, 2026
\end{titlepage}

% Abstract
\section*{Abstract}

Global reanalysis products such as ERA5-Land provide spatially complete weather fields but at resolutions too coarse for local applications, particularly in mountainous regions where temperature can vary by several degrees over short distances. This project investigates Convolutional Conditional Neural Processes (ConvCNPs) for statistical downscaling of daily maximum temperature from the ${\sim}$11\,km ERA5-Land grid to ${\sim}$1\,km resolution over Switzerland, building upon the architecture of \citet{vaughan2022} and adapting it to the topographically complex Swiss domain with high-resolution elevation features from the swisstopo DHM25. The best model, trained on ten years of data (2014--2023) with five-fold temporal cross-validation, achieves a mean absolute error of 1.31\,\textdegree C and a CRPS-based skill score of 0.524 relative to bilinear interpolation, reducing the expected prediction error by more than half. An ablation study reveals that the elevation MLP is the indispensable component---without it, the model diverges entirely---while the Topographic Position Index provides a secondary benefit and explicit seasonal features in the MLP are mildly detrimental. Under sparse on-grid input the model degrades gracefully, maintaining positive skill down to approximately 10\% of the input grid; however, zero-shot deployment on off-grid station observations does not achieve positive skill at any density tested. All configurations exhibit severely overconfident uncertainty estimates, a structural limitation of the Gaussian likelihood training objective. These results demonstrate that ConvCNPs are a viable and effective approach to climate downscaling in complex terrain, and identify uncertainty calibration and native support for non-gridded input as the key challenges for operational deployment.

% Acknowledgements
\section*{Acknowledgements}

\par{
I wish to express my sincere gratitude to my supervisor Dr. Christian Donner, for his invaluable expert guidance and kind support throughout the duration of this project, and to the Swiss Data Science Center, for the opportunity, support and resources to conduct this project. I am also grateful to Ann-Kristin Grimm, Vit Hornacek, Dmitry Fedoruk and Brad Kratochvil, for the support and flexibility to balance my professional responsibilities with my studies. Finally, my deepest thanks go to my wife and son, Isabel and André, for the motivation to pursue this adventure, for their continuous encouragement, and for their patience throughout the demanding periods of my Diploma in Advanced Studies.
}

% Lists
\tableofcontents
\listoffigures
\listoftables

\section*{List of Acronyms}
\begin{description}
    \item[ANP] Attentive Neural Process
    \item[CNN] Convolutional Neural Network
    \item[CNP] Conditional Neural Process
    \item[ConvCNP/CCNP] Convolutional Conditional Neural Process
    \item[ConvNP] Convolutional Neural Process
    \item[CRPS] Continuous Ranked Probability Score (Probabilistic skill score)
    \item[DEM] Digital Elevation Model
    \item[DHM25] Digital Height Model 25\,m (swisstopo)
    \item[ECMWF] European Centre for Medium-Range Weather Forecasts
    \item[ERA5-Land] ECMWF Reanalysis v5 Land (High-resolution terrestrial dataset)
    \item[GP] Gaussian Process
    \item[MAE] Mean Absolute Error
    \item[MLP] Multi-Layer Perceptron
    \item[NP] Neural Process
    \item[NWP] Numerical Weather Prediction
    \item[RMSE] Root Mean Square Error
    \item[TPI] Topographic Position Index
\end{description}

\newpage